\chapter{Dernière diffusion et baryogenèse}

\section{Regarder au loin, c’est regarder dans le passé}

À différentes distances, nous voyons l’Univers à différentes époques. À proximité, le temps rétrospectif n’a pas beaucoup d’importance. L’Univers d’il y a un an, ou d’il y a cent ans, n’était pas très différent de l’Univers actuel. Mais suffisamment loin, nous voyons un Univers plus jeune. Un Univers plus jeune est plus dense, et un Univers plus dense est plus chaud.

À une distance suffisamment grande, l’Univers était à peu près aussi chaud que la surface du Soleil. Nous ne le voyons pas ainsi, car les photons ont été décalés vers le rouge en chemin jusqu’à nous. Un facteur de décalage vers le rouge d’environ mille signifie que la température observée est réduite du même facteur :
\[
T_{\mathrm{obs}} \sim \frac{T_{\mathrm{emit}}}{1+z},
\qquad
1+z \sim 10^3.
\]\footnote{Cette formule est une transcription éditoriale concise de l’argument verbal sur le décalage vers le rouge à 00:02:20--00:02:39.}

C’est là que l’Univers est devenu opaque. Il était principalement constitué d’hydrogène, avec un peu d’hélium. Au-dessus de quelques milliers de degrés, l’hydrogène est ionisé. Au-delà de cette surface, il n’y a donc pas d’atomes neutres, mais des ions : principalement des protons et des électrons libres. Les photons y sont également plus énergétiques, essentiellement des rayons X. Si un atome se forme par hasard, les collisions ou un photon le réionisent rapidement.

Un gaz ionisé est un conducteur électrique, et un conducteur ne laisse pas passer le rayonnement sans le perturber. Les charges réagissent aux champs électriques, oscillent d’avant en arrière et réémettent du rayonnement. Au-delà de cette distance, l’Univers nous est donc opaque. La surface de dernière diffusion ressemble à la surface visible du Soleil : ce n’est pas une paroi matérielle, mais un endroit où le milieu passe assez brusquement de l’opacité à la transparence.

La plupart des particules qui s’y trouvent sont des photons, dans une proportion d’environ \(10^8\). Il en va de même aujourd’hui, bien que les photons soient désormais si peu énergétiques qu’ils ne contribuent pas de façon importante à la densité d’énergie actuelle. Lorsque la température passe sous la température d’ionisation de l’hydrogène, l’Univers devient transparent et nous pouvons voir au travers.

La vision directe s’arrête là. La question suivante est de savoir jusqu’où la physique atomique, nucléaire et des particules que nous connaissons peut encore nous ramener dans le passé.

\section{Remonter plus loin dans le temps}

Juste en deçà de la dernière diffusion, outre les photons, les principales particules sont les protons et les électrons, avec quelques noyaux d’hélium et quelques deutérons. Il n’y a pratiquement ni antiprotons ni positrons. Mais, en remontant plus loin, la température finit par devenir suffisamment élevée pour que deux photons puissent entrer en collision et produire un électron et un positron :
\[
\gamma + \gamma \leftrightarrow e^- + e^+.
\]

Le seuil fondamental est que les photons doivent au moins fournir l’énergie de masse au repos de la paire :
\[
E \ge 2m_e c^2.
\]
En dessous de ce seuil, la lumière ordinaire ne suffit pas. Au-dessus, des paires électron--positron peuvent être créées, et la réaction inverse peut également se produire. Une marmite de matière à cette température contiendrait des réactions entre particules jouant le rôle des réactions chimiques dans la matière ordinaire.

Juste au-dessus du seuil, les photons sont encore plus nombreux que les électrons et les positrons. La masse au repos de l’électron a un coût énergétique élevé. Mais, à mesure que la température dépasse largement le seuil, l’énergie de masse au repos de l’électron devient moins importante que son énergie cinétique, et les densités deviennent comparables. Les électrons, les positrons et les photons ont alors à peu près la même densité et entrent en collision à qui mieux mieux.

\begin{figure}[ht]
\centering
\includegraphics[width=0.62\textwidth]{lecture_06_figure_02.png}
\caption{Schéma du seuil de production de paires comparant les échelles de masse de l’électron et du proton.}
\lectureframe{00:18:47}
\end{figure}

La question suivante est de savoir ce qui arrive aux protons. Ils sont déjà présents et, à ce stade, très nombreux, mais l’énergie thermique ne suffit pas encore à produire en abondance des paires proton--antiproton. La raison est simple :
\[
m_p \approx 2000\,m_e,
\qquad
E \ge 2m_p c^2
\]
le seuil est beaucoup plus élevé.\editorialnote{Dans l’image, l’indication de masse inférieure est moins lisible que l’indication supérieure \(2m_e c^2\), mais la comparaison énoncée porte explicitement sur les seuils de production des électrons et des protons.}

C’est déjà beaucoup plus chaud qu’à l’intérieur d’une étoile, où les antiparticules ne sont pas présentes en abondance notable.

Mais cette description n’est encore que provisoire. Lorsque la température devient assez élevée pour produire librement des protons et des antiprotons, elle l’est aussi pour désintégrer les protons en quarks. Il vaut donc mieux décrire le véritable bouillon brûlant comme un mélange d’électrons, de positrons, de quarks, d’antiquarks, de photons, de neutrinos, de gluons et d’autres espèces. Une réaction schématique simple inscrite au tableau est
\[
e^- + e^+ \to \gamma \to q + \bar q.
\]

\begin{classroomqa}
\lecturetimestamp{00:20:05}
\audiencequestion{S’agit-il de mésons ou de préquarks ?}
\lecturerresponse{Les mésons sont constitués de quarks. Dès que la température est assez élevée pour désintégrer un proton, elle désintègre aussi les mésons.}
\end{classroomqa}

Jusqu’ici, le raisonnement remontait le temps. Inversions maintenant le sens et demandons-nous ce qui se passe lorsque le bouillon chaud et symétrique se refroidit.

\section{Neutralité électrique et Univers chaud symétrique}

À température suffisamment élevée, l’Univers contient d’énormes densités de particules et d’antiparticules. Électrons et positrons apparaissent par paires, tout comme quarks et antiquarks, et l’image naturelle est celle d’abondances presque égales de chaque particule et de son antiparticule.

La neutralité électrique demande quelques précisions. L’Univers est électriquement neutre aujourd’hui. Pour fixer les idées, imaginons non pas notre Univers réel, mais un Univers fermé et fini, comme une hypersphère. Prenons d’abord une sphère à deux dimensions comme modèle. Les lignes de champ électrique ne peuvent commencer ni finir qu’en présence d’une charge. Une charge positive émet des lignes de force ; une charge négative les absorbe. Sur une sphère fermée, ces lignes ne peuvent pas simplement s’échapper à l’infini. Elles convergent de nouveau au point opposé. Cette image aide donc à se représenter qu’un Univers fermé et fini ne peut porter de charge nette. La question de savoir si notre Univers réel est littéralement fermé n’est pas le propos. L’essentiel est que les charges positives et négatives semblent se compenser avec une précision extraordinaire à grande échelle.

La neutralité électrique ne \emph{signifie pas} que le nombre de protons est égal à celui des antiprotons. Elle signifie que la charge positive totale est égale à la charge négative totale. Les protons, les positrons et toutes les autres espèces chargées positivement compensent ensemble les électrons et toutes les autres espèces chargées négativement.

Aussi loin que la physique connue reste fiable, et peut-être un peu au-delà, l’Univers était assez chaud pour que particules et antiparticules soient présentes en abondance. Toute expérience d’accélérateur que l’on puisse citer a eu un analogue dans l’Univers primordial à une certaine époque. L’hypothèse naturelle est celle d’un état initial presque symétrique : autant d’électrons que de positrons, autant de quarks que d’antiquarks et, si l’on prend les protons et les antiprotons comme représentants des baryons et des antibaryons,
\[
N_p = N_{\bar p}.
\]

Cette hypothèse est séduisante parce que les lois de la physique sont presque symétriques entre particules et antiparticules. Un positron se comporte presque, mais pas tout à fait, comme un électron. Cette quasi-symétrie est le problème.

\section{Faire avancer le film}

Supposons que l’Univers commence dans un état chaud, avec des abondances égales ou presque de particules et d’antiparticules, puis qu’il se refroidisse.

Tant que la température est suffisamment élevée, des collisions énergétiques peuvent produire des protons et des antiprotons, qui peuvent s’annihiler et se transformer en photons. Mais la température finit par passer sous le seuil de production des protons. Les paires proton--antiproton peuvent encore s’annihiler en photons, mais, après une certaine expansion, ces photons ont perdu assez d’énergie pour ne plus pouvoir produire à nouveau des paires proton--antiproton. C’est un aller simple. La même chose se produit plus tard pour les électrons et les positrons, à une température bien inférieure au seuil des protons, mais qui reste très supérieure aux températures stellaires.

Comme l’expansion est lente par rapport aux taux des réactions microscopiques, l’Univers reste pour l’essentiel proche de l’équilibre thermique pendant ce processus. C’est ce qu’on entend ici par expansion adiabatique. Plus précisément, l’entropie varie très peu. Pour un bouillon de particules chaudes, l’entropie correspond essentiellement au nombre de particules, à un facteur d’ordre un près. Les protons, antiprotons, électrons et positrons qui s’annihilent sont donc remplacés, à peu près, par des photons.

Si le nombre de protons et d’antiprotons était exactement égal, ils s’annihileraient simplement tous. Il en irait de même pour les électrons et les positrons. Il ne resterait que des photons. Mais, s’il y avait ne serait-ce qu’un très léger excès de protons sur les antiprotons, les paires correspondantes s’annihileraient et l’excédent non apparié subsisterait. Un exemple concret permet de comprendre : s’il y avait cent protons et cinquante antiprotons, les cinquante antiprotons trouveraient chacun un proton avec lequel s’annihiler, et il resterait cinquante protons. La neutralité électrique exige alors un excédent résiduel comparable d’électrons par rapport aux positrons.

Cela permet de déduire l’asymétrie primordiale à partir du nombre actuel de photons. Aujourd’hui, le nombre de photons est environ \(10^9\) à \(10^{10}\) fois supérieur au nombre de protons ou d’électrons :
\[
\frac{N_\gamma}{N_p} \sim 10^9\text{--}10^{10}.
\]
Le bouillon chaud primordial ne contenait donc qu’environ un proton excédentaire par milliard de paires proton--antiproton. Ce minuscule déséquilibre est la raison pour laquelle des atomes existent encore.

\begin{classroomqa}
\lecturetimestamp{00:42:25}
\audiencequestion{Pourquoi ne pas interpréter l’excédent de façon statistique, comme quelque chose qui se serait créé pendant le refroidissement avec une faible probabilité ?}
\lecturerresponse{Pour obtenir un proton excédentaire à partir d’un état initial symétrique, il faudrait créer un proton sans créer aussi un antiproton. Il faudrait également créer avec lui une charge négative, par exemple un électron. Le modèle standard connu ne comporte aucun processus de ce type à basse énergie. Aux énergies testées, le nombre baryonique est conservé ; le refroidissement ordinaire ne produit donc pas d’excédent net.}
\end{classroomqa}

La charge électrique n’est donc pas l’obstacle. Un exemple schématique serait
\[
\gamma + \gamma \to p + e^-.
\]\footnote{Cette formule est une transcription éditoriale de l’exemple verbal donné à 00:44:16--00:46:35. Elle vise à isoler l’idée que la charge électrique pourrait être conservée alors que le nombre baryonique ne le serait pas.}
La charge électrique totale serait alors nulle. Mais une autre règle empirique serait enfreinte. Dans le modèle standard testé, on ne peut pas créer un proton sans créer aussi un antiproton. Pour ce qui nous intéresse ici, la quantité conservée pertinente est le nombre baryonique,
\[
B=(N_p+N_n)-(N_{\bar p}+N_{\bar n}).
\]\footnote{Cette formule concise formalise la définition donnée oralement à 00:45:18--00:45:40.}
Aux énergies connues, cette quantité est conservée ; l’asymétrie ne peut donc pas être produite durant la phase tardive du refroidissement, bien établie expérimentalement.

\section{Baryogenèse}

C’est le sujet qu’on appelle la baryogenèse. Supposons que l’Univers ait réellement commencé avec autant de baryons que d’antibaryons. Que devraient permettre les lois de la physique pour transformer cet état symétrique en le léger excédent observé ?

Au tableau, cela a été ramené à deux expressions symboliques,
\[
N_p=N_{\bar p},
\qquad
N_p-N_{\bar p},
\]
puis à une liste de conditions reliant la première à la seconde.

\begin{figure}[ht]
\centering
\includegraphics[width=0.72\textwidth]{lecture_06_figure_04.png}
\caption{Conditions de production d’une asymétrie baryonique à partir d’un état initial symétrique.}
\lectureframe{01:22:26}
\end{figure}

\editorialnote{Sur cette image, les mots de la liste numérotée ne sont que partiellement lisibles. Les trois éléments ci-dessous reprennent la discussion orale qui l’accompagne.}

La première condition est la violation du nombre baryonique. Si le nombre baryonique est une loi de conservation exacte, un Univers initialement doté d’un nombre baryonique net nul le conserve.

La deuxième condition est un biais entre particules et antiparticules. La violation du nombre baryonique ne suffit pas. Si la nature autorise un processus produisant un proton et son processus symétrique produisant un antiproton avec la même probabilité, le déséquilibre moyen reste nul. Il faut que les taux diffèrent.

Les fluctuations statistiques ordinaires sont également beaucoup trop faibles. Si le nombre total de particules est d’environ \(10^{90}\), photons compris, alors une fluctuation de l’ordre de \(\sqrt{N}\) ne donne qu’environ \(10^{45}\), très loin des quelque \(10^{80}\) protons que contient l’Univers.

La troisième condition est que l’Univers ait été, à un moment donné, suffisamment éloigné de l’équilibre. L’argument fait ici appel au théorème CPT. Si un gaz chaud est resté en équilibre de toute éternité, un film projeté à l’envers est identique à celui projeté à l’endroit. Dans ce cas, la symétrie entre particules et antiparticules est effectivement rétablie par la symétrie exacte qui combine la conjugaison de charge, la parité et l’inversion du temps. Pour qu’un biais réel se fige, l’évolution effective doit différer lorsqu’on la suit dans un sens ou dans l’autre, à l’échelle des temps microscopiques. À une époque suffisamment reculée, l’expansion de l’Univers était assez rapide pour que ce soit le cas.

Les trois conditions sont donc :
\begin{enumerate}
\item violation du nombre baryonique ;
\item violation de la symétrie entre particules et antiparticules, plus précisément violation de CP ;
\item une phase hors équilibre thermique.
\end{enumerate}

\begin{classroomqa}
\lecturetimestamp{01:09:08}
\audiencequestion{En mettant ces conditions en correspondance avec CPT, la condition 2 est-elle bien la violation de CP et la condition 3 concerne-t-elle réellement la direction du temps ?}
\lecturerresponse{Oui. La condition 2 devrait être formulée comme une violation de CP. La condition 3 exprime le fait que l’évolution effective était hors équilibre ; ce n’est ni une loi de conservation distincte ni la même chose que la violation du nombre baryonique.}
\end{classroomqa}

Avec ces trois conditions réunies, même un Univers qui commence avec autant de protons que d’antiprotons peut être biaisé dans un sens ou dans l’autre. On ne sait toujours pas si la nature viole réellement le nombre baryonique et, le cas échéant, comment elle le fait.

\section{Ce que l’on sait et ce que l’on ignore}

La question encore ouverte concerne la violation du nombre baryonique. Nous n’avons jamais observé directement dans la nature de réaction qui modifie
\[
(N_p+N_n)-(N_{\bar p}+N_{\bar n}).
\]
Cela peut simplement signifier que les réactions pertinentes exigent des températures bien plus élevées que celles testées directement en laboratoire, ou qu’elles sont extraordinairement improbables à basse énergie. Mais, pour l’instant, cela reste une question ouverte.

C’est pourquoi la désintégration du proton occupe une place si importante dans cette discussion. Si le nombre baryonique n’est pas exact, le proton devrait en principe être instable. Or il est remarquablement stable. La limite inférieure citée ici est de l’ordre de
\[
\tau_p \gtrsim 10^{34}\ \text{years}.
\]
Elle provient d’énormes expériences menées avec de très grands volumes d’eau équipés de détecteurs de lumière et enfouis profondément sous terre. Elle ne réfute pas la violation du nombre baryonique. Elle signifie seulement que toute violation de ce type doit être extraordinairement supprimée à basse énergie.

Une remarque distincte sur la physique connue intervient également dans la discussion. Entre environ \(100\) et \(250\,\mathrm{GeV}\), les interactions faibles deviennent beaucoup plus semblables aux interactions électromagnétiques. Cela peut avoir son importance dans des modèles détaillés, mais ce n’est pas l’étape spectaculaire de l’argument présent.

La matière noire ne change guère la logique fondamentale. Dans la plupart des théories, elle interagit si faiblement avec la matière ordinaire que, pour cette question, elle joue essentiellement un rôle de spectatrice.

\begin{classroomqa}
\lecturetimestamp{01:24:49}
\audiencequestion{Quel est l’effet de la gravité dans tout cela ?}
\lecturerresponse{Principalement de deux façons. Premièrement, sans expansion cosmique, il n’y aurait pas eu de phase hors équilibre. Deuxièmement, l’évaporation des trous noirs donne une bonne raison de soupçonner que le nombre baryonique ne peut pas être une loi de conservation exacte pour toujours. En dehors de cela, les interactions gravitationnelles ordinaires entre particules sont trop faibles pour avoir beaucoup d’effet ici.}
\end{classroomqa}

La question suivante porte sur la charge électrique. Le même argument concernant les trous noirs menace-t-il aussi la conservation de la charge ? La réponse est non. Un trou noir chargé porte des lignes de flux électrique qui s’étendent vers l’extérieur et continuent de signaler sa charge. Un tel trou noir ne s’évapore pas complètement ; il tend plutôt vers un résidu extrémal. L’argument des trous noirs menace donc le nombre baryonique, mais pas la charge électrique.

Une autre question demande si le bouillon de quarks très dense de l’Univers primordial doit être simplement considéré comme un trou noir. La réponse est non. Il ne s’agit pas d’un amas dense et statique. L’Univers était en expansion. L’histoire des trous noirs est une expérience de pensée sur le devenir du nombre baryonique, et non une description littérale de l’Univers primordial comme région piégée.

\begin{classroomqa}
\lecturetimestamp{01:31:14}
\audiencequestion{Saurait-on comment modifier le modèle standard pour y inclure de telles réactions ?}
\lecturerresponse{Le problème est inverse. Il est facile d’introduire trop de violation du nombre baryonique et de faire désintégrer le proton beaucoup trop vite. La vraie difficulté consiste à préserver la stabilité du proton telle qu’on l’observe tout en autorisant les processus souhaités dans l’Univers primordial.}
\end{classroomqa}

\section{Entropie et dernières questions sur les neutrons}

Revenons maintenant au bilan. Dans un Univers en expansion adiabatique, l’entropie est approximativement conservée. Pour un bouillon chaud de particules élémentaires, l’entropie correspond essentiellement au nombre de particules, à un facteur d’ordre un près. C’est pourquoi la discussion n’a pas besoin de supposer une production tardive massive de photons à partir de rien. Les annihilations redistribuent principalement le contenu du bouillon sous forme de rayonnement.

La formulation ancienne du problème inversait le raisonnement. Au lieu de demander pourquoi si peu de protons ont survécu par rapport aux photons, on posait la question dans l’autre sens : pourquoi y a-t-il tant de photons par rapport aux protons ? Mais la question plus précise n’est pas de savoir pourquoi l’Univers a produit autant d’entropie tardivement. Il s’agit de comprendre pourquoi l’asymétrie entre particules et antiparticules était si faible dès le départ : environ un proton excédentaire seulement par milliard de paires proton--antiproton.

Un neutron libre se désintègre, mais ce processus ne \emph{viole pas} le nombre baryonique :
\[
n \to p + e^- + \bar\nu.
\]
Le neutron porte le même nombre baryonique que le proton.

Pour un neutron libre, l’énergie disponible pour la désintégration est très faible, car la masse du neutron n’est que légèrement supérieure à celle de l’état final. C’est l’une des raisons pour lesquelles sa désintégration est si lente à l’échelle des particules élémentaires. Lorsqu’un neutron est lié dans un noyau, l’énergie de liaison peut absorber cette faible énergie disponible ; le neutron lié peut alors être stable.

\begin{classroomqa}
\lecturetimestamp{01:43:03}
\audiencequestion{Les neutrons disparaissent-ils simplement de l’Univers, ou s’en crée-t-il aussi ?}
\lecturerresponse{La première réponse est que les neutrons libres disparaissent par désintégration. Mais ce n’est pas toute l’histoire. Il existe de véritables sources de neutrons, notamment les réactions nucléaires dans les étoiles et la formation d’étoiles à neutrons. Le nombre total de neutrons ne diminue donc pas simplement de façon continue.}
\end{classroomqa}